\documentclass[11pt]{article}
\usepackage[T1]{fontenc}\usepackage[utf8]{inputenc}
\usepackage{lmodern,amsmath,amssymb,amsthm,mathtools,microtype,xcolor}
\usepackage[margin=1in]{geometry}
\usepackage[colorlinks=true,linkcolor=blue!50!black,urlcolor=blue!50!black,citecolor=blue!50!black]{hyperref}
\newtheorem{theorem}{Theorem}[section]\newtheorem{proposition}[theorem]{Proposition}
\newtheorem{lemma}[theorem]{Lemma}\newtheorem{corollary}[theorem]{Corollary}
\theoremstyle{remark}\newtheorem{remark}[theorem]{Remark}
\newcommand{\E}{\mathbb E}\newcommand{\Prb}{\mathbb P}\newcommand{\R}{\mathbb R}
\newcommand{\normmax}[1]{\lVert#1\rVert_{\max}}
\setlength{\emergencystretch}{3em}
\hypersetup{pdftitle={An explicit sufficient width for mean-field Bayesian circuit learning},pdfauthor={ChatGPT-6 Astra; Kaarel Hanni; Dmitry Vaintrob}}
\title{An explicit sufficient width for\\mean-field Bayesian circuit learning}
\author{ChatGPT-6 Astra \and Kaarel H\"anni \and Dmitry Vaintrob\thanks{Authors are listed in alphabetical order.}}
\date{October 15, 2026}
\begin{document}\maketitle
\begin{abstract}
For fixed-depth tanh networks with mean-field Gaussian scaling, a global
Boolean margin restriction, and exact noiseless Bayesian conditioning, we
quantify the width threshold in a polynomial-sample learning theorem.
The proof replaces an asymptotic large-deviation comparison by a finite-width
method-of-types bound. Gaussian innovations are partitioned into finitely many
cells, and each empirical type is lifted to a finite-entropy Gaussian
conditional law. Entropy controls both the lift's second moment and its tail
mass; these controls propagate through the hidden covariances.
The resulting threshold is explicit and depends on the target only through
its bounded-weight network complexity. It is double exponential in input
dimension, while the sample bound above the threshold remains polynomial in
input dimension and target complexity. A polynomial threshold, and a
polynomial sample bound valid at every width above target complexity, remain
open for deep networks. We include the variational and compression ingredients
so that the result can be read independently.
\end{abstract}

\section{Model and result}
Let $d\geq1$, $L\geq2$, $m=2^d$, and $r=L-1$. The input space is
$\mathcal X=\{-1,1\}^d$, with augmented input $\bar x=(1,x)$.
Every hidden layer has width $N$, and
\[
 h^0(x)=\bar x,\qquad h^s(x)=\tanh(W^sh^{s-1}(x))\quad(1\leq s<L),
 \qquad z_w(x)=W^Lh^{L-1}(x).
\]
There are no additional biases. The independent centered Gaussian weights
have standard deviations $d^{-1/2}$ in the first layer, $N^{-1/2}$ in
internal layers, and $N^{-1}$ in the readout. Their law is $\mu_N$.
For a Boolean function $g$ write
\[
 A_{N,g}=\{w:g(x)z_w(x)>a\text{ for every }x\},\qquad a=1/2,
 \qquad A_N=\bigcup_g A_{N,g}.
\]
The function prior is $\pi_N(g)=\mu_N(A_{N,g})/\mu_N(A_N)$.
Learning conditions this prior on agreement with independent noiseless
examples $(x_i,f(x_i))$, with $x_i\sim D$. We denote the posterior by
$\Pi_N(\cdot\mid S_n)$ and its average error by
\[
 R_D(\Pi_N)=\E_{g\sim\Pi_N}\Prb_{x\sim D}\{g(x)\ne f(x)\}.
\]
Let $C=C_L(f)$ be the least integer for which $f$ has a depth-$L$,
width-$C$ implementation with every weight bounded by $C$ and strict global
margin $1/2$.

Our threshold uses the following explicit target-cost bound. Set
\begin{align*}
 b_L&=2+\sqrt{2\log(2L)},& A_*&=2+3C^2+b_L,\\
 \rho&=[2+9C^2(d+L+1)A_*^{L-2}]^{-1},\\
 B_0&=d(d+1)C^2+2(L-2)C^4+36C^4,\\
 A&=2B_0+4[d+2+(L-2)C]/\rho+2,& K&=4A+4.
\end{align*}
The cylinder construction in Appendix~\ref{app:cloning} proves, for all
$N\geq C$,
\begin{equation}\label{eq:target-mass}
 \mu_N(A_{N,f})\geq e^{-AN/2}.
\end{equation}
All these quantities are polynomially bounded in $d,C$ at fixed $L$.

Define the quantization parameters
\begin{align}
 B&=4K+2m\log2,& b&=2\sqrt{mB},& \beta&=\frac{1}{128mK},\notag\\
 c&=\left(\frac{\beta}{4b^2}\right)^{2^{r-1}},&
 t&=\frac c8,& \Delta&=\frac{c}{8\sqrt{2m}},\notag\\
 R^2&=4\left[\frac{K+\log2}{t}+\frac m2\log2\right],&&
 q&=1+\left\lceil1+\frac{2R\sqrt m}{\Delta}\right\rceil^m.
 \label{eq:quantization-parameters}
\end{align}
Here $q$ is an upper bound on the number of quantization cells; it is an
integer, and should not be confused with a prediction-confidence threshold.
For $\varepsilon\in(0,1)$ set
\begin{equation}\label{eq:width-threshold}
 N_0=\left\lceil\max\left\{C,1,
 \left(\frac{8rq}{K}\right)^2,
 \frac{4m\log2}{K},\frac8K\log\frac6\varepsilon\right\}\right\rceil.
\end{equation}

\begin{theorem}[Explicit sufficient width]\label{thm:main}
At every fixed depth $L$, there is a polynomial $p_L$ such that the following
holds for every target of finite complexity $C=C_L(f)$, every
$\varepsilon,\delta\in(0,1)$, every input distribution $D$, and every
$N\geq N_0$ from \eqref{eq:width-threshold}.
If
\[
 n\geq\frac6\varepsilon\left[p_L(d,C)+\log\frac1\delta\right],
\]
then, with probability at least $1-\delta$ over the training sample,
\[
 R_D(\Pi_N)\leq\frac\varepsilon3,\qquad
 \Prb_{x\sim D}\!\left\{
 \Pi_N(g(x)=f(x)\mid S_n)>\frac23\right\}\geq1-\varepsilon.
\]
An explicit substitute for $p_L(d,C)$ is $B_L(d,2K)$ from the finite-network
count in Appendix~\ref{app:rate-counting}.
The threshold satisfies
\[
 N_0\leq C+\exp\!\left[O_L\bigl(2^d(d+\log K)\bigr)\right]
       +O\!\left(K^{-1}\log\frac1\varepsilon\right).
\]
In particular, the width threshold is quantified, but is not polynomial in
input dimension. The theorem does not establish the same polynomial sample
bound for every width $N\geq C$.
\end{theorem}

The sample bound and the width threshold play distinct roles. The sample
bound is polynomial in the target complexity and independent of $N$.
The threshold is much larger: quantizing all $2^d$ coordinates of a
Gaussian field leads to an exponential number of cells in $2^d$.
For one hidden layer, an existing direct norm argument already gives a
polynomial sample bound for every $N\geq C$; the present threshold is useful
as a general-depth quantitative statement.

\section{An exact representation and a variational complexity}
Enumerate the inputs by $1,\ldots,m$ and let $\Gamma_m=\mathcal N(0,I_m)$.
For a positive semidefinite matrix $Q$, put
\[
 T_Q(u)=\tanh(Q^{1/2}u),\qquad
 F_Q(u)=T_Q(u)T_Q(u)^{\mathsf T},\qquad
 Q_0(x,y)=\bar x\cdot\bar y/d.
\]
Given probability laws $\eta_1,\ldots,\eta_r$ on $\R^m$ and $v\in\R^m$,
recursively set
\[
 Q_s=\int F_{Q_{s-1}}(u)\,d\eta_s(u),\qquad z=Q_r^{1/2}v.
\]
Define
\begin{equation}\label{eq:innovation-rate}
 I(z)=\inf\left\{\sum_{s=1}^rD(\eta_s\Vert\Gamma_m)+\frac12\|v\|_2^2:
          (\eta_1,\ldots,\eta_r,v)\text{ gives }z\right\},
 \qquad
 J(g)=\inf_{g_xz_x\geq a\ \forall x}I(z).
\end{equation}
Entropy contraction makes this equal to the hidden-feature formula proved
in Appendix~\ref{sec:critical-variational}. We use that appendix and the
quadratic-tilt compression in subsequent appendices only to obtain the
polynomial count of $J$ sublevels. The new finite-width probability estimate
below does not invoke an unquantified large-deviation limit.

For each layer $s$ let $U_{s,1},\ldots,U_{s,N}$ be independent standard
Gaussian vectors, independently across layers, and let
$\eta_{s,N}=N^{-1}\sum_i\delta_{U_{s,i}}$. Independently take
$V_N\sim\mathcal N(0,I_m/N)$. The recursion above gives empirical
covariances $Q_{s,N}$. Exactly for every finite $N$,
\begin{equation}\label{eq:exact-output}
 Z_N\stackrel{\mathrm{law}}=Q_{r,N}^{1/2}V_N.
\end{equation}
Indeed, conditional on the preceding layer, the next preactivation rows are
independent Gaussian fields with covariance equal to that layer's empirical
activation covariance. The readout has conditional covariance $Q_{r,N}/N$.
This proves the representation, including singular covariance matrices.

\section{A finite-entropy lift of an empirical type}
Partition the Euclidean ball $\{u:\|u\|_2\leq R\}$ into intersections with
axis-aligned grid cells of side $\Delta/\sqrt m$, and add its complement as
one outside cell. Ignore cells of zero Gaussian probability; the sampled innovations avoid their union almost surely. Every inside
cell has diameter at most $\Delta$, and their total number, including the
outside cell, is at most $q$ in \eqref{eq:quantization-parameters}.
Write the Gaussian cell probabilities as $\alpha_j$.

\begin{lemma}[Entropy controls the lift]\label{lem:lift}
For a cell-frequency vector $p=(p_j)$ define
\[
 \eta^p=\sum_jp_j\,\Gamma_m(\,\cdot\mid\text{cell }j).
\]
Then $D(\eta^p\Vert\Gamma_m)=D(p\Vert\alpha)$.
If this entropy is at most $K$, then
\[
 \int\|u\|_2^2\,d\eta^p(u)\leq B,
 \qquad p_{\mathrm{out}}\leq t.
\]
\end{lemma}
\begin{proof}
The lift has density $p_j/\alpha_j$ on cell $j$, proving the entropy identity.
The entropy inequality at exponent $1/4$ gives
\[
 \frac14\int\|u\|_2^2\,d\eta^p
 \leq D(\eta^p\Vert\Gamma_m)
      +\log\E_{\Gamma_m}e^{\|U\|_2^2/4}
 \leq K+\frac m2\log2.
\]
For an initially unbounded test function this follows by truncation and
monotone convergence. It is the stated moment bound.
Gaussian Markov gives
\[
 \alpha_{\mathrm{out}}\leq2^{m/2}e^{-R^2/4}
 \leq\exp[-(K+\log2)/t].
\]
Data processing to the two-cell inside/outside partition and the bound on
binary entropy imply
\[
 D(p\Vert\alpha)\geq
 p_{\mathrm{out}}\log(1/\alpha_{\mathrm{out}})-\log2.
\]
It follows that $p_{\mathrm{out}}\leq t$.
\end{proof}

\begin{lemma}[Uniform covariance comparison]\label{lem:quantization}
Suppose an empirical measure $\eta$ has cell type $p$, with
$D(p\Vert\alpha)\leq K$. For every positive semidefinite $Q$ with diagonal
entries at most two,
\[
 \left\|\int F_Q\,d\eta-\int F_Q\,d\eta^p\right\|_{\max}\leq c.
\]
For any two positive semidefinite matrices $Q,Q'$ and any lift of entropy
at most $K$,
\[
 \left\|\int F_Q\,d\eta^p-\int F_{Q'}\,d\eta^p\right\|_{\max}
 \leq b\sqrt{\|Q-Q'\|_{\max}}.
\]
\end{lemma}
\begin{proof}
Each entry of $F_Q$ is bounded in absolute value by one and is
$2\|Q^{1/2}\|_{\mathrm{op}}$-Lipschitz. If the diagonal of $Q$ is bounded
by two, then $\|Q\|_{\mathrm{op}}\leq\operatorname{tr}Q\leq2m$.
Couple each empirical point with an independent conditional Gaussian draw
in its own cell. Inside cells the difference is at most
$2\sqrt{2m}\Delta$; outside it is at most two. Lemma~\ref{lem:lift}
therefore bounds the integrated difference by
\[
 2\sqrt{2m}\Delta+2t=c/2\leq c.
\]
For the second claim, the positive square-root inequality gives
\[
 \|Q^{1/2}-(Q')^{1/2}\|_{\mathrm{op}}
 \leq\|Q-Q'\|_{\mathrm{op}}^{1/2}
 \leq\sqrt{m\|Q-Q'\|_{\max}}.
\]
The entrywise product bound and Cauchy--Schwarz, with
Lemma~\ref{lem:lift}, give an upper bound of
$2\sqrt{m\|Q-Q'\|_{\max}}\sqrt B$, as required.
\end{proof}

\section{A nonasymptotic rate-sublevel bound}
\begin{proposition}[Low type cost implies low variational complexity]
\label{prop:deterministic}
Let $p_s$ be the types of any empirical innovation measures
$\eta_{s,N}$, and let $v\in\R^m$. If
\[
 H:=\sum_{s=1}^rD(p_s\Vert\alpha)+\frac12\|v\|_2^2\leq K,
\]
and $Q_{r,N}^{1/2}v$ has strict margin $a$ and sign $g$, then
$J(g)\leq16K/9<2K$.
\end{proposition}
\begin{proof}
Replace each empirical law by its conditional Gaussian lift $\eta_s^*$
and let $Q_s^*$ be the resulting recursive covariances.
Both recursions start at $Q_0$, whose diagonal is at most two; every later
covariance has diagonal at most one. For
$e_s=\|Q_{s,N}-Q_s^*\|_{\max}$, Lemma~\ref{lem:quantization} gives
\[
 e_0=0,\qquad e_s\leq c+b\sqrt{e_{s-1}}.
\]
Since $b\geq1$ and $c\leq1$, induction yields
\[
 e_s\leq4b^2c^{1/2^{s-1}},\qquad e_r\leq\beta.
\]
For the induction step, $b\sqrt{4b^2c^{1/2^{s-2}}}
=2b^2c^{1/2^{s-1}}$, and the additional $c$ is bounded by the same
power of $c$.
The lifted output $z^*=(Q_r^*)^{1/2}v$ satisfies
\[
 \|z^*-Q_{r,N}^{1/2}v\|_\infty
 \leq\sqrt{m\beta}\|v\|_2
 \leq\sqrt{2mK\beta}=a/4.
\]
Thus $(4/3)z^*$ has strict margin $a$ and sign $g$.
It is obtained from the lifted laws and output vector $(4/3)v$.
Its cost in \eqref{eq:innovation-rate} is at most
\[
 \sum_sD(p_s\Vert\alpha)+\frac{16}{9}\frac12\|v\|_2^2
 \leq\frac{16}{9}H\leq\frac{16}{9}K.
\]
This proves the claim.
\end{proof}

\begin{theorem}[Finite-width aggregate bound]\label{thm:aggregate}
Let $\mathcal A_K=\{g:J(g)\leq2K\}$. At every width $N\geq1$,
\begin{equation}\label{eq:aggregate}
 \mu_N\left(\bigcup_{g\notin\mathcal A_K}A_{N,g}\right)
 \leq (N+1)^{rq}\,2^{m/2}e^{-NK/2}.
\end{equation}
\end{theorem}
\begin{proof}
The probability of a fixed empirical type $p$ from $N$ independent draws
with cell probabilities $\alpha$ is at most $e^{-ND(p\Vert\alpha)}$.
Indeed, the multinomial coefficient is at most $e^{NH(p)}$.
There are at most $(N+1)^q$ possible types per layer.
For a fixed tuple of types put $D_* =\sum_sD(p_s\Vert\alpha)$.
Independence across layers and from $V_N$ gives
\begin{align*}
 &\Prb\!\left\{\text{these types},\ D_*+\tfrac12\|V_N\|_2^2>K\right\}\\
 &\quad\leq e^{-ND_*}
       \exp\{ND_*/2-NK/2\}
       \E\exp\{N\|V_N\|_2^2/4\}\\
 &\quad=2^{m/2}\exp\{-ND_*/2-NK/2\}
 \leq2^{m/2}e^{-NK/2}.
\end{align*}
The exponential bound remains valid when $D_*>K$.
By Proposition~\ref{prop:deterministic}, every full-margin draw whose sign
function lies outside $\mathcal A_K$ has total cost greater than $K$.
Summing over at most $(N+1)^{rq}$ tuples proves \eqref{eq:aggregate}.
\end{proof}

This controls the entire union of undesired functions. It does not sum an
asymptotic estimate separately over all $2^{2^d}$ Boolean functions. The
large dimensional dependence enters instead through quantizing Gaussian
innovations in $\R^{2^d}$.

\section{Proof of the learning theorem and the remaining goal}
\begin{proof}[Proof of Theorem~\ref{thm:main}]
The target-mass bound, Theorem~\ref{thm:aggregate}, and cancellation of
the global-margin normalization give
\[
 \frac{\pi_N(\mathcal A_K^c)}{\pi_N(f)}
 \leq\exp\left[rq\log(N+1)+\frac m2\log2
                   -\frac N2(K-A)\right].
\]
For $N\geq1$, $\log(N+1)\leq\sqrt N$. The first two nonconstant
constraints in \eqref{eq:width-threshold} therefore give
\[
 rq\log(N+1)\leq NK/8,\qquad (m/2)\log2\leq NK/8.
\]
Since $A<K/4$, the ratio is at most $e^{-NK/8}\leq\varepsilon/6$.
In particular $f\in\mathcal A_K$: otherwise the ratio would be at least
one. Appendix~\ref{app:rate-counting} gives
$\log|\mathcal A_K|\leq B_L(d,2K)$, a polynomially bounded quantity.

Every particular function with $D$-error greater than $\varepsilon/6$
survives $n$ noiseless examples with probability at most
$e^{-n\varepsilon/6}$. A union bound over $\mathcal A_K$ shows that,
except on a sample event of probability at most $\delta$, every surviving
function in this set has error at most $\varepsilon/6$ for the stated $n$.
For every sample the posterior denominator is at least $\pi_N(f)$;
therefore its mass outside $\mathcal A_K$ is at most $\varepsilon/6$.
The total posterior error is at most $\varepsilon/3$.
On each input where the correct-label probability is at most $2/3$, the
posterior error probability is at least $1/3$. The mass of these inputs is
at most $3R_D(\Pi_N)\leq\varepsilon$. This proves the strict confidence
claim, including equality in the failure event.

Finally, at fixed $L$, every reciprocal precision and truncation radius in
\eqref{eq:quantization-parameters} is polynomially bounded in $m,K$.
Hence $\log q=O_L(m\log(mK))$. Substitution into
\eqref{eq:width-threshold} proves the asserted growth bound.
\end{proof}

The next research goal is a width threshold polynomial in $d,C$ and the
accuracy parameters, followed ideally by a polynomial sample guarantee at
every $N\geq C$. The present proof does not supply either improvement.
Its costly step is a uniform quantization of all coordinates of a Gaussian
field on the entire cube. The quadratic-tilt compression used for counting
is dimension efficient, but it applies to variational minimizers. A
finite-width counterpart that controls the aggregate mass of networks
without first quantizing an arbitrary $2^d$-dimensional field would be needed
to obtain a substantially sharper threshold by this strategy.

\appendix
For compatibility with the variational notation below, write $\phi=\tanh$ and $\mu_{N,1}=\mu_N$.
% Standalone section: macros deliberately limited to standard AMS commands.
\section{An all-width target-mass bound}\label{app:cloning}
The following construction bounds the target's prior cost at every width $N\geq C$,
at the mean-field exponent.

\begin{proposition}[Uniform cloning mass bound]
Fix $L\geq2$ and a bounded $1$-Lipschitz activation. There is a function $Q_L(d,C)$ bounded by a polynomial such that for every Boolean target with
$C_L(f)\leq C$, and every $N\geq C$,
\[
 -\frac{2}{N}\log\Pr\{f(x)z_w(x)>1/2\text{ for all }x\}
 \leq Q_L(d,C).
\]
The prior variances and absence of later-layer biases are exactly those in the
main model.
\end{proposition}
\begin{proof}
Choose a width-$C$ teacher, increase its readout by a factor of three, and write
its hidden matrices as $B^s$ and readout as $v$. Thus the teacher has margin
strictly greater than $3/2$, hidden entries bounded by $C$, and readout entries
bounded by $3C$. If the original witness has smaller width, first use its actual
width in the construction below; all estimates only improve.
Partition $\{1,\ldots,N\}$ into $C$ groups of sizes $k_a$ differing by at most
one. Then
\[
 1\leq k_a,\qquad \frac{\max k_a}{\min k_a}\leq2,
 \qquad \frac N{k_a}\leq2C.
\]
Let $u_a=k_a^{-1/2}{\bf1}_{\text{group }a}$ and let $P$ be the orthogonal
projection onto their span. The exact clone $w_*$ has first-layer row $B^1_b$
for a neuron in group $b$, hidden entry $B^s_{ba}/k_a$ from group $a$ to group
$b$, and readout $v_b/k_b$ on group $b$. Every hidden activation of the exact
clone is group-constant. Moreover,
\[
 \|W_*^s\|_{\mathrm{op}}\leq\sqrt2 C^2\quad(2\leq s<L),
 \qquad \sqrt N\|W_*^L\|_2\leq3\sqrt2 C^2.
\]
The first bound follows by restricting to the group-constant subspace: the
resulting matrix has entries $\sqrt{k_b/k_a}B^s_{ba}$, and its Frobenius norm
is at most $\sqrt2 C^2$.

Set $t=1$. The random coordinates
$\sqrt d W^1_{ij}$, $\sqrt N W^s_{ij}$ for $2\leq s<L$, and
$\alpha_j=NW^L_j$ are independent centered Gaussians of standard deviation
$t$. Impose the following cylinder event $\mathcal C_\rho$, for
$0<\rho\leq1$:
\begin{itemize}
\item every $\sqrt d W^1_{ij}$ is within $\rho t$ of
$\sqrt d (W_*^1)_{ij}$;
\item for every hidden row $j$ and group $a$, the coordinate
$\langle\sqrt N W^s_{j,\cdot},u_a\rangle$ is within $\rho t$ of its
cloned value $\sqrt{N/k_a}B^s_{ba}$;
\item every $\alpha_j$ is within $\rho t$ of its cloned value $Nv_b/k_b$.
\end{itemize}
All these constraints concern independent scalar Gaussian coordinates.
The orthogonal coordinates of each internal row remain independent prior
Gaussians. Consequently write
\[
 W^s=W_*^s+R^s+E^s,\qquad R^s=R^sP,\quad E^sP=0.
\]
On the cylinder, $\|R^s\|_{\mathrm{op}}\leq\|R^s\|_F\leq\rho t\sqrt C$.
The free matrix is distributed as $(t/\sqrt N)G^s(I-P)$, with $G^s$ a standard
Gaussian matrix, independently of the cylinder coordinates. Let
\[
 b_L=2+\sqrt{2\log(2L)}.
\]
The Gaussian operator-norm bound and a union bound imply, conditionally on
$\mathcal C_\rho$, that all $\|E^s\|_{\mathrm{op}}\leq b_Lt$ with probability
at least $1/2$. This is also valid when $L=2$, where there are no such matrices.

Let $e_s(x)=N^{-1/2}\|h^s_w(x)-h^s_*(x)\|_2$. The input constraints give
$e_1(x)\leq\rho(d+1)/\sqrt d\leq\rho(d+1)$. Since $E^sh^{s-1}_*=0$,
boundedness and Lipschitz continuity give, simultaneously for every input,
\[
 e_s(x)\leq (\sqrt2 C^2+\sqrt C+b_L)e_{s-1}(x)+\rho\sqrt C.
\]
Put $A=2+3C^2+b_L$. Induction yields
$e_s(x)\leq\rho(d+L+1)A^{L-2}$ for $s<L$. At the output,
\[
 |z_w(x)-z_*(x)|
 \leq3\sqrt2 C^2 e_{L-1}(x)+\rho.
\]
Hence choosing
\[
 \rho=\bigl[2+9C^2(d+L+1)A^{L-2}\bigr]^{-1}
\]
ensures output discrepancy at most $1/2$. All networks in the cylinder with
the operator-norm event therefore implement $f$ with margin greater than one.

It remains to price the cylinder. Its number of constrained coordinates is
\[
 m=N\bigl[d+2+(L-2)C\bigr].
\]
The sum of their squared centers is at most $NB_0$, where one may take
\[
 B_0=d(d+1)C^2+2(L-2)C^4+36C^4.
\]
Indeed, a first-layer center has magnitude at most $\sqrt d C$, an internal
center at most $\sqrt{2C}C$, and a readout center at most $6C^2$.
For $Z\sim N(0,t^2)$ and $|Z-a|\leq\rho t$,
\[
 \Pr\{|Z-a|\leq\rho t\}
 \geq\frac{2\rho}{\sqrt{2\pi}}
       \exp\{-a^2/t^2-\rho^2\}.
\]
Multiplication and $t=1$ show
\[
 -\log\Pr(\mathcal C_\rho)
 \leq N B_0
 +N[d+2+(L-2)C]\log\frac{e\sqrt{2\pi}}{2\rho}.
\]
The desired probability is at least $\Pr(\mathcal C_\rho)/2$. Since $N\geq1$, this proves the claim with
\[
 Q_L(d,C)=2B_0+\frac{4[d+2+(L-2)C]}{\rho}+2.
\]
Here we used $\log(e\sqrt{2\pi}/(2\rho))\leq2/\rho$ for $0<\rho\leq1$.
The displayed expression is a polynomial in $d,C$ at every fixed depth.
\end{proof}

\section{An exact variational principle at the critical exponent}
\label{sec:critical-variational}

For arbitrary fixed depth, the critical prior admits a direct finite-domain
large-deviation description. This result requires neither a propagation-of-chaos
assumption nor differentiability of the activation. It provides the variational
foundation for the compression argument that follows.

Fix $\gamma=1$, enumerate $\mathcal X=\{x^1,\ldots,x^m\}$, where $m=2^d$,
and put $r=L-1$. All limits in this section hold with $d$ and $L$ fixed.
Write $E=[-1,1]^m$ and let $\mathcal P(E)$ carry its weak topology.
For $\nu\in\mathcal P(E)$ define
\[
 K(\nu)=\int_E hh^{\mathsf T}\,d\nu(h),\qquad
 Q_0(a,b)=\frac{\bar x^a\cdot\bar x^b}{d}.
\]
For any positive semidefinite $m\times m$ matrix $Q$, let
\[
 \mathsf T_Q(u)=\bigl(\phi((Q^{1/2}u)_a)\bigr)_{a=1}^m,
 \qquad
 \mathsf M_Q=(\mathsf T_Q)_\#\mathcal N(0,I_m).
\]
Here $T_\#\rho$ is the distribution of $T(U)$ when $U$ has distribution $\rho$,
and $Q^{1/2}$ is the positive semidefinite square root. Relative entropy is
\[
 D(\nu\Vert\rho)=
 \begin{cases}
 \displaystyle\int\log\frac{d\nu}{d\rho}\,d\nu,&\nu\ll\rho,\\
 +\infty,&\text{otherwise}.
 \end{cases}
\]
Finally define the Gaussian energy
\[
 \mathcal G_Q(z)=
 \begin{cases}
 \frac12 z^{\mathsf T}Q^+z,&z\in\operatorname{range}(Q),\\
 +\infty,&z\notin\operatorname{range}(Q),
 \end{cases}
\]
where $Q^+$ is the Moore--Penrose inverse.

\begin{theorem}[Critical finite-domain large deviations]
\label{thm:critical-ldp}
Let $Z_N=(z_w(x^a))_{a=1}^m$ under the unconditioned prior
$\mu_{N,1}$. The sequence $Z_N$ satisfies a large-deviation principle
with speed $N$ and good rate function
\begin{equation}\label{eq:critical-rate}
 \begin{split}
 I(z)=\inf_{\nu_1,\ldots,\nu_r\in\mathcal P(E)}
 \bigg\{&D(\nu_1\Vert\mathsf M_{Q_0})
  +\sum_{s=2}^{r}D\bigl(\nu_s\Vert\mathsf M_{K(\nu_{s-1})}\bigr)
  +\mathcal G_{K(\nu_r)}(z)\bigg\}.
 \end{split}
\end{equation}
In particular, for every open $O\subseteq\mathbb R^m$ and closed
$F\subseteq\mathbb R^m$,
\[
 -\inf_{z\in O}I(z)
 \leq\liminf_{N\to\infty}\frac1N\log\Prb(Z_N\in O),
 \qquad
 \limsup_{N\to\infty}\frac1N\log\Prb(Z_N\in F)
 \leq-\inf_{z\in F}I(z).
\]
The infimum in \eqref{eq:critical-rate} is attained whenever it is finite.
No covariance matrix in the statement is required to be nonsingular.
\end{theorem}

\begin{proof}
We give an exact representation of the finite-width law and then apply
Sanov's theorem and continuous contraction~\cite{dembozeitouni}.

Let $\Gamma_m=\mathcal N(0,I_m)$. Independently for each
$s=1,\ldots,r$, let $U_{s,1},\ldots,U_{s,N}$ have law $\Gamma_m$, and set
\[
 \eta_{s,N}=\frac1N\sum_{i=1}^N\delta_{U_{s,i}}.
\]
Also let $V_N=G/\sqrt N$, where $G\sim\Gamma_m$ is independent of these
empirical measures. Recursively define
\[
 \nu_{1,N}=(\mathsf T_{Q_0})_\#\eta_{1,N},\qquad
 \nu_{s,N}=(\mathsf T_{K(\nu_{s-1,N})})_\#\eta_{s,N}\quad(2\leq s\leq r).
\]
These are exactly the joint laws of the network's empirical hidden-activation
vectors on $\mathcal X$. Indeed, first-layer preactivation vectors are
independent Gaussians with covariance $Q_0$. Conditional on the preceding
layer, each subsequent row of weights gives an independent Gaussian vector
with covariance
\[
 \frac1N\sum_{i=1}^N
 \bigl(h_i^{s-1}(x^a)h_i^{s-1}(x^b)\bigr)_{a,b}
 =K(\nu_{s-1,N}).
\]
The conditional distribution depends on the previous layer only through its
empirical measure. Finally, readout weights have variance $N^{-2}$, so
\begin{equation}\label{eq:critical-output-representation}
 Z_N\ \stackrel{\mathrm{law}}{=}\ K(\nu_{r,N})^{1/2}V_N.
\end{equation}
This representation is exact for every $N$.

Sanov's theorem on $\mathcal P(\mathbb R^m)$ gives the good rate
$D(\eta\Vert\Gamma_m)$ for each $\eta_{s,N}$. The elementary Gaussian
large-deviation principle gives rate $\|v\|_2^2/2$ for $V_N$.
Independence and the finite-product principle therefore give the good rate
\[
 \sum_{s=1}^r D(\eta_s\Vert\Gamma_m)+\frac12\|v\|_2^2
\]
for $(\eta_{1,N},\ldots,\eta_{r,N},V_N)$.

The recursive map from these variables to $(\nu_1,\ldots,\nu_r,z)$ is
continuous. To verify the only potentially delicate point, suppose
$Q_j\to Q$ and $\eta_j\Rightarrow\eta$ weakly. For every continuous
$b:E\to\mathbb R$, the functions $b\circ\mathsf T_{Q_j}$ converge
uniformly on each compact subset of $\mathbb R^m$ to
$b\circ\mathsf T_Q$, since matrix square roots are continuous and $\phi$
is continuous. Their absolute values are uniformly bounded. Tightness of a
weakly convergent sequence of probability measures then gives
\[
 \int b\circ\mathsf T_{Q_j}\,d\eta_j
 \longrightarrow\int b\circ\mathsf T_Q\,d\eta.
\]
Also $\nu\mapsto K(\nu)$ is continuous because every coordinate
$h_ah_b$ is continuous and bounded on $E$. This proves continuity of the
whole recursive map, including \eqref{eq:critical-output-representation}.

Continuous contraction now gives the rate by minimizing over its preimages.
For a measurable map $T$ and a reference probability $\rho$, the entropy
contraction identity is
\begin{equation}\label{eq:entropy-contraction-identity}
 \inf_{\eta:T_\#\eta=\nu}D(\eta\Vert\rho)
 =D(\nu\Vert T_\#\rho).
\end{equation}
The lower bound is data processing. When the right side is finite and
$q=d\nu/d(T_\#\rho)$, the measure $d\eta=q(T(u))\,d\rho(u)$ attains
equality. If the right side is infinite, data processing already proves
the identity in the extended real sense. Apply
\eqref{eq:entropy-contraction-identity} successively with
$T=\mathsf T_{Q_0}$ and $T=\mathsf T_{K(\nu_{s-1})}$.
The last minimization is
\[
 \inf_{v:Q^{1/2}v=z}\frac12\|v\|_2^2=\mathcal G_Q(z),
\]
including singular $Q$. This proves \eqref{eq:critical-rate}.

The contraction theorem supplies a good rate and exponential tightness.
For completeness, the displayed minimization is also directly well behaved:
relative entropy is jointly lower semicontinuous, the map
$\nu\mapsto\mathsf M_{K(\nu)}$ is weakly continuous, and
\[
 \mathcal G_Q(z)=\sup_{t\in\mathbb R^m}
 \left\{t^{\mathsf T}z-\frac12t^{\mathsf T}Qt\right\}
\]
is jointly lower semicontinuous in $(Q,z)$. Since $\mathcal P(E)^r$ is
compact, every finite infimum in \eqref{eq:critical-rate} is attained.
Moreover $K(\nu)\preceq mI_m$, so $I(z)\geq\|z\|_2^2/(2m)$.
\end{proof}

\begin{corollary}[Full-margin masses and posterior rates]
\label{cor:critical-margin-rates}
For a Boolean function $f$, regarded as a vector in $\{-1,1\}^m$, set
\[
 J(f)=\inf_{z:\ f_a z_a\geq1/2\ \forall a}I(z).
\]
Then the strict-margin event used throughout this paper has the exact rate
\begin{equation}\label{eq:critical-full-margin-rate}
 \lim_{N\to\infty}-\frac1N\log\mu_{N,1}(A_{N,f})=J(f).
\end{equation}
For a fixed noiseless training set $S$, let $\mathcal B_S$ be the Boolean
functions agreeing with its labels. If
$J_S=\min_{g\in\mathcal B_S}J(g)<\infty$, the actual globally conditioned
posterior satisfies, for every $f\in\mathcal B_S$,
\[
 \lim_{N\to\infty}-\frac1N\log\Pi_{N,1}(f\mid S)=J(f)-J_S.
\]
In particular, posterior mass concentrates on the minimizers of $J$ within
$\mathcal B_S$. These logarithmic asymptotics do not determine how mass is
divided between tied minimizers.
\end{corollary}

\begin{proof}
For $c\geq1$, retaining any feature laws in \eqref{eq:critical-rate} and
scaling only the output gives
\[
 I(cz)\leq c^2 I(z),
\]
because the entropy terms are nonnegative and the Gaussian energy scales
quadratically. If $f_a z_a\geq1/2$ for all $a$, then $cz$ has strict
margin for every $c>1$. Letting $c\downarrow1$ shows that the infimum of
$I$ over the open strict-margin orthant equals its infimum over the
closed orthant. The lower and upper large-deviation bounds therefore agree,
proving \eqref{eq:critical-full-margin-rate}.

Conditioning the prior on the global margin event and subsequently on the
labels gives exactly
\[
 \Pi_{N,1}(f\mid S)=
 \frac{\mu_{N,1}(A_{N,f})}
 {\sum_{g\in\mathcal B_S}\mu_{N,1}(A_{N,g})}.
\]
For fixed $d$ this is a finite sum, whose exponential rate is the smallest
$J(g)$ among its terms. This proves the ratio formula and concentration
claim.
\end{proof}

The variational principle exhibits the distinction between the critical
regime and a frozen kernel model. At $\gamma=1$, deviations of a hidden
feature law have an order-$N$ entropy cost, at the same scale as an order-one
output. Thus the optimization is over every hidden feature law, with a
separate entropy charge for each layer. It is not in general legitimate to
set those laws to their unconditioned values. The following sections obtain a
polynomial description of the minimizing feature laws for $\tanh$ and use it
to compress every Boolean function with bounded full-margin rate.

\section{Quadratic tilts at a critical full-margin minimizer}
\label{sec:critical-stationarity}

We now specialize to $\phi=\tanh$. Its analyticity and its bounded first two
derivatives yield additional structure at the minimizers of
\eqref{eq:critical-rate}. The following argument uses the exact variational
principle above; it does not assume mean-field equations.

\begin{lemma}[Positive definite feature covariances]
\label{lem:critical-positive-definite}
If $(\nu_1,\ldots,\nu_r,z)$ has finite cost in
\eqref{eq:critical-rate}, then every $K(\nu_s)$ is positive definite.
\end{lemma}
\begin{proof}
For any nonzero $c\in\mathbb R^m$, the function
\[
 H_c(b,v)=\sum_{a=1}^m c_a\tanh(b+v\cdot x^a)
\]
is real analytic and is not identically zero. To see the latter, choose
$v_0$ such that the numbers $v_0\cdot x^a$ are all distinct. If $H_c$ were
identically zero, evaluating it at $(b,v)=(-t\theta,tv_0)$ and letting
$t\to\infty$ would give
\[
 \sum_a c_a\operatorname{sign}(v_0\cdot x^a-\theta)=0
\]
for every $\theta$ different from these numbers. The difference between
two successive such identities is twice one coefficient $c_a$, so all
coefficients would vanish. A nonzero real-analytic function has a
Lebesgue-null zero set. The first-layer weights have a positive Gaussian
density on the $(d+1)$-dimensional $(b,v)$ space. Consequently
\[
 \mathsf M_{Q_0}\{h:c^{\mathsf T}h=0\}=0.
\]
Finite entropy implies $\nu_1\ll\mathsf M_{Q_0}$, and therefore
$c^{\mathsf T}K(\nu_1)c>0$ for every $c\ne0$.

Inductively, if $Q=K(\nu_{s-1})$ is positive definite, then
$\mathsf M_Q$ has a positive density on $(-1,1)^m$, since coordinatewise
$\tanh$ is a smooth bijection from $\mathbb R^m$ onto that set.
Every proper linear hyperplane is $\mathsf M_Q$-null. Absolute continuity
$\nu_s\ll\mathsf M_Q$ again gives positive definiteness of $K(\nu_s)$.
\end{proof}

\begin{theorem}[Quadratic form of the optimal feature laws]
\label{thm:critical-quadratic-tilts}\label{thm:critical-tilts}
Let $a>0$, $f\in\{-1,1\}^m$, and suppose
\[
 J_a(f)=\inf_{z:\ f_bz_b\geq a\ \forall b}I(z)\leq K<\infty.
\]
There is a minimizing tuple $(\nu_1,\ldots,\nu_r,z)$.
Write $Q_s=K(\nu_s)$ and $\beta=Q_r^{-1}z$.
There are symmetric matrices $\Gamma_1,\ldots,\Gamma_r$ such that
\begin{equation}\label{eq:critical-quadratic-tilts}
 \frac{d\nu_1}{d\mathsf M_{Q_0}}(h)
   =\frac{e^{h^{\mathsf T}\Gamma_1h}}{Z_1},\qquad
 \frac{d\nu_s}{d\mathsf M_{Q_{s-1}}}(h)
   =\frac{e^{h^{\mathsf T}\Gamma_sh}}{Z_s}\quad(2\leq s\leq r).
\end{equation}
Moreover,
\begin{equation}\label{eq:critical-beta-bound}
 \|\beta\|_1\leq\frac{2K}{a},\qquad
 \Gamma_r=\frac12\beta\beta^{\mathsf T}.
\end{equation}
If $g_s=\sum_{i,j}|(\Gamma_s)_{ij}|$, then
\begin{equation}\label{eq:critical-gamma-recursion}
 g_r\leq\frac{2K^2}{a^2},\qquad
 g_s\leq3g_{s+1}+2g_{s+1}^2\quad(1\leq s<r).
\end{equation}
Thus, at fixed depth, all these entrywise norms are bounded by a polynomial
in $K$ and $1/a$, independently of the number $m=2^d$ of inputs.
\end{theorem}

\begin{proof}
Existence follows from the goodness of $I$, closedness of the margin
constraints, and attainment of the feature-law infimum in
Theorem~\ref{thm:critical-ldp}. Every $Q_s$ is positive definite by
Lemma~\ref{lem:critical-positive-definite}.

For fixed feature laws, the output minimizes the strictly convex quadratic
$z^{\mathsf T}Q_r^{-1}z/2$ subject to the linear margin constraints.
Their feasible set has a strictly feasible point. The elementary
quadratic-program optimality conditions give multipliers $\lambda_b\geq0$
such that
\[
 \beta_b=f_b\lambda_b,\qquad
 \lambda_b(f_bz_b-a)=0.
\]
It follows that
\[
 \frac12\beta^{\mathsf T}z
 =\frac a2\sum_b\lambda_b
 =\frac a2\|\beta\|_1\leq J_a(f)\leq K,
\]
which proves the first bound in \eqref{eq:critical-beta-bound}.

We justify the feature-law variations, including possible zeros of their
densities. Fix a layer $s$, freeze every other variable, and denote its
reference measure by $\mu_s$ (either $\mathsf M_{Q_0}$ or
$\mathsf M_{Q_{s-1}}$). Its contribution to the objective has the form
\[
 D(\nu_s\Vert\mu_s)+F_s(K(\nu_s))+\text{constant},
\]
where $F_s$ is differentiable near the current positive definite $Q_s$.
For $s=r$, $F_r(Q)=z^{\mathsf T}Q^{-1}z/2$, so
\[
 \nabla F_r(Q_r)=-\frac12\beta\beta^{\mathsf T}.
\]
For $s<r$, put $\xi(h)=(\operatorname{arctanh}(h_b))_{b=1}^m$ and
$B_{s+1}=\int\xi(h)\xi(h)^{\mathsf T}\,d\nu_{s+1}(h)$.
This moment is finite: by invariance of relative entropy under the
coordinatewise $\tanh$ bijection, the preactivation law has finite
relative entropy with respect to $\mathcal N(0,Q_s)$, and the Gaussian
exponential moment of a sufficiently small multiple of
$\xi^{\mathsf T}Q_s^{-1}\xi$ controls its expectation by the entropy
inequality. The Gaussian density ratio, whose Jacobian cancels under
$\tanh$, gives, near $Q_s$,
\[
 F_s(Q)=\frac12\log\det Q+\frac12\operatorname{tr}(Q^{-1}B_{s+1})
       +\text{constant}.
\]
In particular,
\begin{equation}\label{eq:critical-Gamma-from-moments}
 -\nabla F_s(Q_s)
 =\frac12Q_s^{-1}(B_{s+1}-Q_s)Q_s^{-1}=:\Gamma_s.
\end{equation}

Put $B=\nabla F_s(Q_s)$ and let $\rho_s$ have density
$e^{-h^{\mathsf T}Bh}/Z$ with respect to $\mu_s$. Its logarithm is bounded,
so $D(\nu_s\Vert\rho_s)$ is finite. For
$\nu_{s,t}=(1-t)\nu_s+t\rho_s$, convexity of relative entropy gives
\[
 D(\nu_{s,t}\Vert\rho_s)-D(\nu_s\Vert\rho_s)
 \leq-tD(\nu_s\Vert\rho_s).
\]
The identity
$D(\nu\Vert\rho_s)=D(\nu\Vert\mu_s)+\int h^{\mathsf T}Bh\,d\nu+\log Z$
and differentiability of $F_s$ imply that the full objective changes by
at most $-tD(\nu_s\Vert\rho_s)+o(t)$. Minimality forces
$D(\nu_s\Vert\rho_s)=0$. Thus $\nu_s=\rho_s$, proving
\eqref{eq:critical-quadratic-tilts} with
$\Gamma_s=-\nabla F_s(Q_s)$. This also proves
$\Gamma_r=\beta\beta^{\mathsf T}/2$.

It remains to bound the apparently ill-conditioned expression
\eqref{eq:critical-Gamma-from-moments} without using $\|Q_s^{-1}\|$.
Under $\nu_{s+1}$ the preactivation vector $\xi$ has Gaussian reference
$\mathcal N(0,Q_s)$ tilted by $e^{V(\xi)}$, where
\[
 V(\xi)=\tanh(\xi)^{\mathsf T}\Gamma_{s+1}\tanh(\xi).
\]
Twice applying Gaussian integration by parts gives
\[
 B_{s+1}-Q_s
 =Q_s\,\mathbb E_{\mathrm{tilt}}
       [\nabla^2V+(\nabla V)(\nabla V)^{\mathsf T}]\,Q_s.
\]
This identity is justified directly because $V$ and its first two
derivatives are bounded. Substitution into
\eqref{eq:critical-Gamma-from-moments} gives
\[
 \Gamma_s=\frac12\mathbb E_{\mathrm{tilt}}
       [\nabla^2V+(\nabla V)(\nabla V)^{\mathsf T}].
\]
The bounds $|\tanh|\leq1$, $|\tanh'|\leq1$, and
$|\tanh''|\leq2$ imply
\[
 \|\nabla V\|_1\leq2g_{s+1},\qquad
 \sum_{i,j}|\partial_i\partial_jV|\leq6g_{s+1}.
\]
Therefore $g_s\leq3g_{s+1}+2g_{s+1}^2$.
The final-layer bound is
$g_r=\|\beta\|_1^2/2\leq2K^2/a^2$.
Iterating the displayed quadratic recursion a fixed number of times proves
the asserted polynomial dependence.
\end{proof}

\section{Compression of bounded quadratic-tilt towers}\label{sec:critical-compression}
This section uses $\phi=\tanh$. Write $\|A\|_{\max}=\max_{ij}|A_{ij}|$ and
$\|A\|_{1,\mathrm{ent}}=\sum_{ij}|A_{ij}|$ for matrices on the finite input
set. These are different norms; their dual pairing is what avoids a factor
of $2^d$ below.
For a positive semidefinite covariance $Q$ and a symmetric matrix $\Gamma$,
let
\[
 \mathsf F_\Gamma(Q)=
 \frac{\mathbb E_{Z\sim\mathcal N(0,Q)}
 [\phi(Z)\phi(Z)^\top\exp\{\phi(Z)^\top\Gamma\phi(Z)\}]}
 {\mathbb E_{Z\sim\mathcal N(0,Q)}
 [\exp\{\phi(Z)^\top\Gamma\phi(Z)\}]}.
\]
All activations are coordinatewise. The definition is valid for singular $Q$.

\begin{lemma}[A dimension-independent covariance comparison]\label{lem:tilted-covariance}
If $g=\|\Gamma\|_{1,\mathrm{ent}}$, then for all positive semidefinite
$Q,Q'$ of the same size,
\[
 \|\mathsf F_\Gamma(Q)-\mathsf F_\Gamma(Q')\|_{\max}
 \leq 8(1+g)^2\|Q-Q'\|_{\max}.
\]
\end{lemma}
\begin{proof}
Put $V(z)=\phi(z)^\top\Gamma\phi(z)$ and
$Q_t=(1-t)Q+tQ'$. For a twice differentiable bounded test function $F$, let
$\mathbb E_t$ denote expectation under the Gaussian of covariance $Q_t$
reweighted by $e^V$. Gaussian differentiation and the quotient rule give
\begin{align*}
 \frac{d}{dt}\mathbb E_tF
 =\frac12\sum_{a,b}(Q'-Q)_{ab}\bigl\{
 &\mathbb E_t[\partial_{ab}F+
       \partial_aF\partial_bV+\partial_bF\partial_aV]\\
 &+\operatorname{Cov}_t(F,\partial_{ab}V+
                                   \partial_aV\partial_bV)\bigr\}.
\end{align*}
This formula has no reciprocal partition-function bound: all terms are
expectations under a probability measure. Define
$M_1(H)=\sup_z\sum_a|\partial_aH(z)|$ and
$M_2(H)=\sup_z\sum_{a,b}|\partial_{ab}H(z)|$.
Because $|\phi|,|\phi'|\leq1$ and $|\phi''|\leq2$,
\[
 M_1(V)\leq2g,\qquad M_2(V)\leq6g.
\]
For $F(z)=\phi(z_i)\phi(z_j)$, including $i=j$,
$\|F\|_\infty\leq1$, $M_1(F)\leq2$, and $M_2(F)\leq6$.
The absolute derivative is therefore at most
\[
 \frac12\|Q-Q'\|_{\max}
 \{6+8g+2(6g+4g^2)\}
 =(3+10g+4g^2)\|Q-Q'\|_{\max}.
\]
Integrating in $t$ proves the assertion, with the displayed slightly larger
constant. One may initially replace both covariances by $Q+\epsilon I$ and
$Q'+\epsilon I$. The functions and all derivatives used above are bounded,
and the partition functions lie between $e^{-g}$ and $e^g$, so weak Gaussian
convergence lets $\epsilon\downarrow0$. This also proves the singular case.
\end{proof}

\begin{proposition}[Finite realization of a bounded-tilt tower]\label{prop:tilt-compress}
Fix $L\geq2$, $K\geq1$, and $G\geq1$. Let $Q_0$ be the covariance of the
Gaussian first-layer field $w\cdot\bar x$, with
$w\sim\mathcal N(0,I_{d+1}/d)$, indexed by $x\in\mathcal X$.
Suppose symmetric matrices $\Gamma_1,\ldots,\Gamma_{L-1}$ and a vector
$\beta\in\mathbb R^{\mathcal X}$ satisfy
\[
 \|\Gamma_s\|_{1,\mathrm{ent}}\leq G,\qquad
 \|\beta\|_1\leq4K,
 \qquad Q_s=\mathsf F_{\Gamma_s}(Q_{s-1}).
\]
Then the function $z(x)=(Q_{L-1}\beta)_x$ can be uniformly approximated to
accuracy $1/8$ by a depth-$L$ network whose width and weights are bounded by
$\operatorname{poly}_L(d,K,G)$.
If $f(x)z(x)\geq1/2$ everywhere, a factor-two increase in the finite
network's readout implements $f$ with margin at least $3/4$.
\end{proposition}
\begin{proof}
Let $b=8(1+G)^2$ and
\[
 \eta=[32K(L-1)b^{L-2}]^{-1}.
\]
Choose the common width $n$ to be the least integer satisfying
\[
 n\geq2\eta^{-2}\log\bigl(8(L-1)\,2^{2d}\bigr).
\]
The construction is random only to prove existence.
The matrices $\Gamma_s$ may be indexed by the entire cube; they select the
sampling laws, but are not stored or evaluated by the resulting finite
network. No efficient sampling or construction algorithm is claimed.
At the first layer, independently sample $n$ Gaussian weight vectors
$w_i\sim\mathcal N(0,I_{d+1}/d)$ tilted by
\[
 \exp\{h(w_i)^\top\Gamma_1h(w_i)\},\qquad
 h(w_i)_x=\phi(w_i\cdot\bar x).
\]
Their activation vectors have exactly the law defining
$Q_1=\mathsf F_{\Gamma_1}(Q_0)$.
At each subsequent hidden layer, conditional on the already sampled previous
activation vectors $h_1,\ldots,h_n$, independently sample $n$ standard
Gaussian row vectors $\xi_j\in\mathbb R^n$ tilted by
\[
 \exp\{\phi(Z_j)^\top\Gamma_s\phi(Z_j)\},
 \qquad Z_j(x)=\frac1{\sqrt n}\sum_{i=1}^n\xi_{ji}h_i(x).
\]
Use the actual network weights $\xi_{ji}/\sqrt n$.
If $\widehat Q_{s-1}=n^{-1}\sum_i h_ih_i^\top$ is the preceding empirical
activation covariance, then, conditionally, $Z_j$ before tilting is Gaussian
with exactly that covariance. Thus the new activation vectors are
conditionally independent, with second moment
$\mathsf F_{\Gamma_s}(\widehat Q_{s-1})$.

Every covariance entry is the empirical mean of $n$ variables in $[-1,1]$.
Hoeffding's inequality, a union bound over $2^{2d}$ entries, and conditioning
at successive layers show, with probability at least $3/4$, that
\[
 \|\widehat Q_s-
          \mathsf F_{\Gamma_s}(\widehat Q_{s-1})\|_{\max}\leq\eta
 \quad(1\leq s<L),
\]
where $\widehat Q_0=Q_0$.
Lemma~\ref{lem:tilted-covariance} then implies
\[
 \|\widehat Q_{L-1}-Q_{L-1}\|_{\max}
 \leq (L-1)b^{L-2}\eta.
\]
For each last-layer activation vector $h_i$, choose the deterministic
readout weight $v_i=(\beta\cdot h_i)/n$. The resulting output is
$\widehat z(x)=(\widehat Q_{L-1}\beta)_x$, and therefore
\[
 \|\widehat z-z\|_\infty
 \leq4K(L-1)b^{L-2}\eta=1/8.
\]
Although $\beta$ is indexed by the entire input cube, its only role in the
finite network is to select these $n$ real readout weights; each satisfies
$|v_i|\leq4K/n$. It is not an additional computational input to the network.

It remains to bound the sampled weights without losing polynomial size.
For every sampled row, its tilted density relative to its Gaussian reference
is at most $e^{2G}$, since the quadratic tilt lies in $[-G,G]$.
Let $D_n=n(d+1)+(L-2)n^2$ be the total number of incoming scalar coordinates,
and choose
\[
 B=\bigl[4G+2\log(8\max\{D_n,1\})\bigr]^{1/2}.
\]
Each raw incoming coordinate has reference variance at most one, so, also
conditionally on all preceding layers, its probability of exceeding $B$ in
absolute value is at most $2e^{2G-B^2/2}$.
A union bound gives probability at most $1/4$ that any incoming coordinate
exceeds $B$. The actual internal weights divide the raw coordinates by
$\sqrt n$ and so satisfy the same bound. Intersecting this event with the
covariance-approximation event leaves probability at least $1/2$.
Fix a realization in the intersection. Both $n$ and
$\max\{B,8K/n\}$ after the factor-two readout increase are bounded by
polynomials in $d,K,G$ at fixed depth. The stated margin follows from
$2(1/2-1/8)=3/4$.
\end{proof}

\begin{lemma}[Counting finite realizations]\label{lem:tilt-count}
Fix $L$, width $n$, and an entrywise bound $B\geq1$. The Boolean functions
having depth-$L$, width-$n$ implementations with margin at least $3/4$ and
weights bounded by $B$ number at most
\[
 \exp\{P\log(3+16B\Lambda)\},\qquad
 P=n(d+1)+(L-2)n^2+n,
 \quad \Lambda=(d+Ln)[1+n(B+1)]^{L-1}.
\]
\end{lemma}
\begin{proof}
Round every weight to a multiple of $(8\Lambda)^{-1}$.
The first hidden layer's coordinatewise error is at most $(d+1)/(8\Lambda)$.
Every subsequent hidden layer has error at most $n(B+1)$ times the preceding
error plus $n/(8\Lambda)$; the same estimate applies at the linear readout.
Induction bounds the final error by $1/8$. Each rounded coordinate has at
most $3+16B\Lambda$ possible values, so the $P$-parameter grid has the stated
size. Rounding preserves every represented Boolean sign function.
\end{proof}

\section{Critical rate complexity and Bayesian generalization}\label{app:rate-counting}
\begin{theorem}[Polynomial rate-sublevel counting]\label{thm:critical-counting}
For every fixed depth $L\geq2$, with activation $\tanh$, there is a polynomial
$B_L$ such that, for every $K\geq1$,
\[
 \log\#\{f:\mathcal X\to\{-1,1\}:J(f)\leq K\}\leq B_L(d,K).
\]
Moreover every function in this sublevel set has a depth-$L$ implementation
of width and weight magnitude at most $\operatorname{poly}_L(d,K)$.
\end{theorem}
\begin{proof}
Take a minimizing tuple for $J(f)\leq K$ in
Theorem~\ref{thm:critical-tilts}, with margin $a=1/2$.
Its vector $\beta$ satisfies $\|\beta\|_1\leq4K$.
Let $g_r^*=8K^2$ and recursively
$g_s^*=3g_{s+1}^*+2(g_{s+1}^*)^2$ for $s<r=L-1$.
Then $G=1+\max_s g_s^*$ is a polynomial in $K$ at fixed depth,
and all the quadratic tilt matrices have entrywise norm at most $G$.
The feature covariances satisfy
$Q_s=\mathsf F_{\Gamma_s}(Q_{s-1})$, and the optimal output is
$z=Q_r\beta$.
Proposition~\ref{prop:tilt-compress} supplies a finite implementation with
margin at least $3/4$, common width bounded by a polynomial in $d,K$, and
weights bounded by another such polynomial.
The bounds in its construction can be chosen uniformly for every function
in the sublevel set. Lemma~\ref{lem:tilt-count} counts the resulting fixed
parameter grid and gives the stated bound.
For the complexity statement, increase width by zero padding if necessary
so that the width also bounds every weight magnitude. This preserves the
network for $\tanh$, since $\tanh(0)=0$.
\end{proof}



\paragraph{An explicit count.}
For completeness, a numerical expression for $B_L(d,U)$, $U\geq1$, follows
directly from the preceding proofs. Put $g_r=8U^2$,
$g_s=3g_{s+1}+2g_{s+1}^2$, $G=1+\max_s g_s$, and
\[
 b_*=8(1+G)^2,\qquad
 \eta_*=[32Ur b_*^{r-1}]^{-1},\qquad
 n_* =\left\lceil2\eta_*^{-2}\log(8r\,2^{2d})\right\rceil.
\]
Let $D_*=n_*(d+1)+(L-2)n_*^2$ and
\[
 B_* =\max\{1,\sqrt{4G+2\log(8\max\{D_*,1\})},8U/n_*\},
 \qquad P_*=D_*+n_*.
\]
\[
 \Lambda_*=(d+Ln_*)[1+n_*(B_*+1)]^{L-1}.
\]
Then one may take
\[
 B_L(d,U)=P_*\log(3+16B_*\Lambda_*).
\]
It is polynomially bounded at each fixed depth. It is a count of Boolean
functions, not a claim that the optimal tilts can be computed efficiently.

\begin{thebibliography}{9}
\bibitem{dembozeitouni} Amir Dembo and Ofer Zeitouni.
\emph{Large Deviations Techniques and Applications}. Second edition.
Springer, 1998.
\bibitem{companion} ChatGPT-6 Astra, Kaarel H\"anni, and Dmitry Vaintrob.
\emph{Bayesian circuit learning at mean-field scaling: a variational principle
and width-independent generalization}. Companion manuscript, October 2026.
The variational, stationarity, and compression proofs reproduced in the
appendices originate in that manuscript.
\end{thebibliography}
\end{document}
